%% Hyperparameter sensitivity. Numbers from code/results/bench/{ra_mu*,ra_ab_*,ra_omega*} (100 portfolios).
\subsection{\blt{Hyperparameter Sensitivity}} \label{s:sensitivity}
Tables~\ref{tab:sensitivity} and \ref{tab:sensitivity_paired} vary, one at a time, the three user-set quantities that the $2^k$ design tuned and the bound $\omega_{\max}$ that Sec.~\ref{s:results} identified as the effective risk weight; 100 portfolios, everything else as in Table~\ref{tab:hyper}.

\emph{Prior shrinkage $\mu$.} Smaller $\mu$ means a weaker prior (worth $\mu/(1-\mu)$ times the 125 days it is built from) and, through the marginal likelihood, a sparser structure: the genetic algorithm keeps 15 of 137 regressors at $\mu=0.7$ against 57 at $\mu=0.9$ and 78 at $\mu=0.95$. $\mu=0.7$ improves the Sharpe ratio by $0.18$ ($p<0.001$, better in 79\% of the portfolios) and reduces the turnover from $15\%$ to $11\%$ per day; $\mu=0.95$ is worse. The value $0.9$ chosen on 2021--2022 data is therefore not the best value on 2023--2025, but the ranking against the benchmarks does not change: even $\mu=0.7$ stays below the equal-weight portfolio.

\emph{Forgetting prior $(\alpha_0,\beta_0)$.} The optimised forgetting weights follow their prior closely (median $\gamma$ of $0.68$, $0.57$, $0.12$ and $0.00$ for $\gamma_0=0.6$, $0.41$, $0.1$ and $0.05$), but none of the three alternatives changes any performance metric by more than $0.01$ in Sharpe ratio. Together with the ablation of Sec.~\ref{s:ablation_results} this shows that the data-based forgetting is inert on this data: what the model predicts is fixed by the prior and the structure, not by how the statistic is updated.

\emph{The risk weight.} Capping the search interval of $\omega^*$ exposes the mechanism of Sec.~\ref{s:results}. With $\omega_{\max}=1$ the residual-minimising search of the implementation returns the lower end of the interval ($\omega\approx10^{-5}$), the return term disappears from the loss, the optimal action is $a_t=a_{t-1}$, and the policy is the rebalanced equal-weight portfolio to three decimals. With $\omega_{\max}=10$ or $100$ (the search returns the bound) the return forecasts enter the loss with a weight that lets them dominate the variance term, the policy takes corner positions in single stocks, turns the portfolio over 1.6 and 0.9 times per day, and loses $73\%$ and $46\%$ of the capital, as the risk-neutral ablation does. The same happens with the three principled elicitations of $\omega$ we tried in place of the search bound (the $\omega$ that makes the KLD between the model and the ideal equal to $1/2$, the $\omega$ that makes the ideal's mean shift equal to the predicted return, and a fixed $\omega=40$; their $\omega$ lies between 27 and 91 in 80\% of the decision steps): they lose $63$--$65\%$ of the capital with a daily turnover of $1.3$. Only $\omega=10^3$, where the variance term dominates by two orders of magnitude and the forecasts are effectively ignored, gives a usable policy. The reason is the forecasts themselves. On the 500 test days of the first ten portfolios the one-step-ahead predictions $\hat y_t$ have a standard deviation of $0.0267$ against $0.0269$ for the realised returns, a correlation with them of $-0.01$ (per asset: $-0.01$), a sign accuracy of $49.7\%$, and a mean squared error twice that of predicting zero (\texttt{bench/forecast\_quality.py}). They are as large as the returns and carry no information, so any loss that takes them at face value trades on noise, and a large $\omega$ is the only setting in which the design does not. The elicitation of $\omega$ in Sec.~\ref{s:FPD}, whatever its analytical form, cannot repair this, and the practical conclusion is that the point forecast must be shrunk (or the design must integrate over the parameter uncertainty that the posterior provides) before it enters the loss.

\begin{table*}[ht!]
\centering
\footnotesize
\setlength{\tabcolsep}{4pt}
\begin{tabular}{lrrrrrrrr}
\hline
strategy & total return & ann. return & ann. vol. & Sharpe (ann.) & max DD & Calmar & turnover/day & cost paid \\
\hline
Risk-Averse (RA) & 1.126 & 5.6\% & 19.4\% & 0.37 & -25.5\% & 0.36 & 14.8\% & 14.8\% \\
RA, $\mu=0.5$ & 1.150 & 7.0\% & 18.4\% & 0.45 & -22.0\% & 0.42 & 16.9\% & 16.9\% \\
RA, $\mu=0.7$ & 1.190 & 8.8\% & 18.4\% & 0.55 & -21.5\% & 0.53 & 10.9\% & 10.9\% \\
RA, $\mu=0.95$ & 1.096 & 4.2\% & 19.7\% & 0.30 & -26.5\% & 0.25 & 18.4\% & 18.4\% \\
RA, $(\alpha_0,\beta_0)=(0.2,0.2)$ & 1.127 & 5.6\% & 19.4\% & 0.37 & -25.5\% & 0.36 & 14.8\% & 14.8\% \\
RA, $(\alpha_0,\beta_0)=(0.3,0.6)$ & 1.122 & 5.5\% & 19.0\% & 0.37 & -24.8\% & 0.37 & 14.9\% & 14.9\% \\
RA, $(\alpha_0,\beta_0)=(0.5,0.45)$ & 1.119 & 5.4\% & 18.8\% & 0.37 & -24.5\% & 0.35 & 15.2\% & 15.2\% \\
RA, $\omega_{\max}=1$ & 1.305 & 14.0\% & 21.0\% & 0.72 & -26.3\% & 0.57 & 1.6\% & 1.6\% \\
RA, $\omega_{\max}=10$ & 0.266 & -50.5\% & 38.7\% & -1.72 & -79.1\% & -0.63 & 159.7\% & 159.7\% \\
RA, $\omega_{\max}=100$ & 0.541 & -27.5\% & 23.3\% & -1.31 & -55.4\% & -0.48 & 92.4\% & 92.4\% \\
RA, $\omega$: unit KLD & 0.353 & -41.8\% & 29.5\% & -1.76 & -70.3\% & -0.58 & 132.2\% & 132.2\% \\
RA, $\omega$: match prediction & 0.367 & -40.4\% & 27.4\% & -1.82 & -68.6\% & -0.58 & 127.8\% & 127.8\% \\
RA, $\omega=40$ & 0.358 & -41.3\% & 28.4\% & -1.80 & -69.7\% & -0.58 & 131.5\% & 131.5\% \\
\hline
\end{tabular}
\caption{Sensitivity to the user-set hyperparameters (prior shrinkage $\mu$, prior forgetting weights) (mean over 100 portfolios).}
\label{tab:sensitivity}
\end{table*}

\begin{table*}[ht!]
\centering
\footnotesize
\setlength{\tabcolsep}{4pt}
\begin{tabular}{lrrrrrrrr}
\hline
benchmark & $\Delta$Sharpe & 95\% CI & $p$ & win rate & $\Delta$max DD & $p$ & $\Delta$total ret. & $p$ \\
\hline
RA, $\mu=0.5$ & -0.08 & [-0.13, -0.03] & 0.005 & 40\% & -3.5\% & $<$0.001 & -0.024 & 0.030 \\
RA, $\mu=0.7$ & -0.18 & [-0.23, -0.13] & $<$0.001 & 21\% & -3.9\% & $<$0.001 & -0.063 & $<$0.001 \\
RA, $\mu=0.95$ & +0.07 & [+0.02, +0.12] & 0.005 & 65\% & +1.1\% & 0.006 & +0.030 & 0.010 \\
RA, $(\alpha_0,\beta_0)=(0.2,0.2)$ & -0.00 & [-0.00, -0.00] & $<$0.001 & 36\% & +0.0\% & 0.201 & -0.001 & $<$0.001 \\
RA, $(\alpha_0,\beta_0)=(0.3,0.6)$ & -0.00 & [-0.02, +0.01] & 0.003 & 68\% & -0.7\% & 0.341 & +0.004 & 0.003 \\
RA, $(\alpha_0,\beta_0)=(0.5,0.45)$ & +0.00 & [-0.02, +0.02] & 0.570 & 53\% & -0.9\% & $<$0.001 & +0.007 & 0.196 \\
RA, $\omega_{\max}=1$ & -0.35 & [-0.44, -0.26] & $<$0.001 & 23\% & +0.9\% & 0.236 & -0.179 & $<$0.001 \\
RA, $\omega_{\max}=10$ & +2.09 & [+1.95, +2.23] & $<$0.001 & 99\% & +53.6\% & $<$0.001 & +0.860 & $<$0.001 \\
RA, $\omega_{\max}=100$ & +1.68 & [+1.60, +1.76] & $<$0.001 & 100\% & +29.9\% & $<$0.001 & +0.585 & $<$0.001 \\
RA, $\omega$: unit KLD & +2.13 & [+2.02, +2.24] & $<$0.001 & 99\% & +44.9\% & $<$0.001 & +0.773 & $<$0.001 \\
RA, $\omega$: match prediction & +2.19 & [+2.08, +2.30] & $<$0.001 & 100\% & +43.1\% & $<$0.001 & +0.759 & $<$0.001 \\
RA, $\omega=40$ & +2.17 & [+2.07, +2.28] & $<$0.001 & 100\% & +44.2\% & $<$0.001 & +0.768 & $<$0.001 \\
\hline
\end{tabular}
\caption{Paired comparison of the default hyperparameters against the alternatives.}
\label{tab:sensitivity_paired}
\end{table*}

